\section{Results}
\label{sec:results}

\subsection{Study 1: Baseline Overconfidence in Equity Markets}
\label{sec:results-study1}

Table~\ref{tab:study1} summarizes hit rates and ECE for the 1-day equity prediction
task ($n = 100$ per model). All three models are substantially overconfident: even the
best-performing model (Claude) achieves only 58\% coverage at the 80\% confidence
level, and only 73\% at 90\%. GPT-4 is dramatically worse---its 80\% CI contains
the true outcome only 31\% of the time, and its 90\% CI only 38\% (ECE = 43.7\%).
These results confirm universal overconfidence and that GPT-4 is the worst-calibrated
model in the simplest possible setting.

\begin{table}[ht]
  \centering
  \caption{Study 1: Overall Calibration in 1-Day Equity Predictions ($n = 100$ per model).
    ECE = Expected Calibration Error; lower is better. Perfect calibration = 0\%.}
  \label{tab:study1}
  \begin{tabular}{lcccc}
    \toprule
    Model   & Hit@50\% & Hit@80\% & Hit@90\% & ECE    \\
    \midrule
    Claude  & 30.0\%   & 58.0\%   & 73.0\%   & 19.7\% \\
    Gemini  & 29.0\%   & 58.0\%   & 65.0\%   & 22.7\% \\
    GPT-4   & 20.0\%   & 31.0\%   & 38.0\%   & 43.7\% \\
    \midrule
    Perfect & 50.0\%   & 80.0\%   & 90.0\%   & 0.0\%  \\
    \bottomrule
  \end{tabular}
\end{table}

\subsection{Study 2: Calibration Across Prediction Horizons}
\label{sec:results-study2}

Table~\ref{tab:study2} shows how calibration changes with prediction horizon for each
model. The results reveal a pronounced and divergent horizon effect. Claude's 80\% hit
rate drops from 80.5\% at 1 day to 46.2\% at 18 days (ECE rises from 1.4\% to 29.8\%),
a pattern consistent with systematic interval-width anchoring on short-horizon
volatility. GPT-4---already poorly calibrated at 1 day (ECE = 21.8\%)---deteriorates
further, reaching 27.2\% hit@80 at 18 days and an ECE of 47.4\%.

Gemini's pattern is strikingly different. Its 1-day calibration (ECE = 1.8\%) is the
best of any model at any horizon. While it degrades at intermediate horizons
(ECE = 16.6\% at 18 days), it recovers at longer windows, achieving 80.1\% hit@80
and ECE = 2.5\% at 22 days---near-perfect calibration. This suggests Gemini's CI
widths scale approximately proportionally to horizon uncertainty in a way that matches
realized volatility at longer windows, even as Claude and GPT-4 fail to do so.

\begin{table}[ht]
  \centering
  \caption{Study 2: 80\% Hit Rate and ECE by Model and Prediction Horizon.
    Hit@80 = proportion of predictions where true outcome falls within 80\% CI.
    ECE = mean absolute calibration error across 50/80/90\% levels.}
  \label{tab:study2}
  % TODO: Add remaining horizon rows (2d, 3d, 7d, 20d, 21d) from study2/data/results/summary.csv
  \resizebox{\columnwidth}{!}{%
  \begin{tabular}{lcccccc}
    \toprule
    Horizon & Cl.\ H80 & Cl.\ ECE & Gem.\ H80 & Gem.\ ECE & GPT H80 & GPT ECE \\
    \midrule
    1d  & 80.5\% & 1.4\%  & 78.5\% & 1.8\%  & 55.2\% & 21.8\% \\
    6d  & 56.0\% & 20.6\% & 71.7\% & 8.4\%  & 34.5\% & 41.6\% \\
    18d & 46.2\% & 29.8\% & 59.8\% & 16.6\% & 27.2\% & 47.4\% \\
    22d & 51.2\% & 26.9\% & 80.1\% & 2.5\%  & 24.4\% & 49.9\% \\
    \bottomrule
  \end{tabular}}
\end{table}

\subsection{Study 3: Domain-Specific Calibration}
\label{sec:results-study3}

Table~\ref{tab:study3-overall} summarizes overall calibration in the multi-domain
study ($n = 68$ per model). Claude exhibits the best overall calibration (ECE = 5.3\%),
followed by Gemini (7.2\%) and GPT-4 (18.9\%), consistent with Study 1 and Study 2
rankings. The overall ECE values are substantially lower than in Studies 1 and 2,
driven by near-perfect calibration in the weather domain (which accounts for 44\% of
Study 3 observations).

\begin{table}[ht]
  \centering
  \caption{Study 3: Overall Calibration by Model ($n = 68$ per model).
    ECE = Expected Calibration Error; lower values indicate better calibration.}
  \label{tab:study3-overall}
  \begin{tabular}{lcccc}
    \toprule
    Model   & Hit@50\% & Hit@80\% & Hit@90\% & ECE    \\
    \midrule
    Claude  & 58.8\%   & 82.4\%   & 85.3\%   & 5.3\%  \\
    Gemini  & 45.6\%   & 73.5\%   & 79.4\%   & 7.2\%  \\
    GPT-4   & 30.9\%   & 63.2\%   & 69.1\%   & 18.9\% \\
    \midrule
    Perfect & 50.0\%   & 80.0\%   & 90.0\%   & 0.0\%  \\
    \bottomrule
  \end{tabular}
\end{table}

Table~\ref{tab:study3-domain} disaggregates calibration by domain. The central result
is a systematic directional reversal: models are overconfident in commodity and equity
markets and underconfident in cryptocurrency and, to a lesser extent, forex. No model
is uniformly overconfident or underconfident, refuting a simple single-mechanism
account of LLM miscalibration.

Commodities showed the most extreme and uniform overconfidence: ECE = 40\% for all
three models, with 80\% and 90\% CIs achieving only 40\% coverage. The cryptocurrency
domain showed a directional inversion: both Claude and GPT-4 achieved 100\% coverage
at the 90\% level, indicating systematic underconfidence. Weather was the
best-calibrated domain overall. GPT-4 performed dramatically worse on forex
(ECE = 44.2\%).

\begin{table}[ht]
  \centering
  \caption{Study 3: Hit Rates and ECE by Model and Domain.
    NBA games excluded (ESPN API). Commod.\ = commodity futures.}
  \label{tab:study3-domain}
  \small
  \begin{tabular}{llccccc}
    \toprule
    Domain     & Model  & $N$ & H50\% & H80\%  & H90\%  & ECE    \\
    \midrule
    Stocks     & Claude & 40  & 50\%  & 85\%   & 90\%   & 1.7\%  \\
               & Gemini & 40  & 25\%  & 70\%   & 75\%   & 16.7\% \\
               & GPT-4  & 40  & 15\%  & 55\%   & 65\%   & 28.3\% \\
    Commod.    & Claude & 10  & 20\%  & 40\%   & 40\%   & 40.0\% \\
               & Gemini & 10  & 20\%  & 40\%   & 40\%   & 40.0\% \\
               & GPT-4  & 10  & 20\%  & 40\%   & 40\%   & 40.0\% \\
    Crypto     & Claude & 20  & 80\%  & 100\%  & 100\%  & 20.0\% \\
               & Gemini & 20  & 20\%  & 80\%   & 100\%  & 13.3\% \\
               & GPT-4  & 20  & 20\%  & 100\%  & 100\%  & 20.0\% \\
    Forex      & Claude & 16  & 38\%  & 75\%   & 75\%   & 10.8\% \\
               & Gemini & 16  & 50\%  & 63\%   & 75\%   & 10.8\% \\
               & GPT-4  & 16  & 25\%  & 25\%   & 38\%   & 44.2\% \\
    Weather    & Claude & 60  & 73\%  & 87\%   & 90\%   & 10.0\% \\
               & Gemini & 60  & 67\%  & 83\%   & 87\%   & 7.8\%  \\
               & GPT-4  & 60  & 47\%  & 77\%   & 80\%   & 5.5\%  \\
    NBA        & All    & 13  & --    & --     & --     & --     \\
    \bottomrule
  \end{tabular}
\end{table}

\subsection{Meta-Knowledge Analysis (Study 3)}
\label{sec:results-metaknowledge}

Table~\ref{tab:mead} presents the Soll-Klayman meta-knowledge ratio for each
model-domain combination. Overall, GPT-4 is the most overconfident ($\mu = 0.81$)
and Claude slightly underconfident in aggregate ($\mu = 1.23$). However, the
aggregate values conceal dramatic variation that constitutes the central finding
of this study.

\begin{table}[ht]
  \centering
  \caption{Study 3: Meta-Knowledge Ratio $\mu = \text{MEAD}/\text{MAD}$ by Model and Domain.
    $\mu < 1$ = overconfidence; $\mu = 1$ = perfect; $\mu > 1$ = underconfidence.}
  \label{tab:mead}
  \begin{tabular}{lccc}
    \toprule
    Domain       & Claude & Gemini & GPT-4 \\
    \midrule
    Stocks       & 0.81   & 0.55   & 0.41  \\
    Commodities  & 0.25   & 0.25   & 0.25  \\
    Crypto       & 14.40  & 1.35   & 1.84  \\
    Forex        & 2.68   & 1.67   & 1.29  \\
    Weather      & 1.28   & 1.04   & 0.85  \\
    \midrule
    Overall      & 1.23   & 0.99   & 0.81  \\
    \bottomrule
  \end{tabular}
\end{table}

All models show $\mu = 0.25$ on commodities: stated uncertainty is only one-quarter
of what actual errors require, the most extreme overconfidence in the dataset. At the
other extreme, Claude's crypto $\mu = 14.4$ indicates CIs 14 times wider than
necessary. The forex domain reveals a compound failure in GPT-4: its high ECE
(44.2\%) is not reflected in its $\mu$ (1.29), because $\mu$ cannot distinguish
misanchored point estimates from miscalibrated CI widths. The two metrics must
be examined jointly to diagnose the full failure mode.
